\documentclass{situationist-preprint}

\title{The Primacy of Perception, Applied to Itself}
\msid{PHE-14-a}
\msfamily{PHE-14}
\msdiscipline{Phenomenology}
\mstype{Principal manuscript}
\mskeywords{primacy of perception; Merleau-Ponty; Husserl; founding; categorial intuition; formal reasoning; evidence; reading and writing}

\begin{document}
\maketitle

\begin{abstract}
Merleau-Ponty held that perception is primary: every form of consciousness, including reflection, science and logic, presupposes the perceptual opening onto the world. The thesis faces its hardest test in formal reasoning, whose objects are not perceived and whose validity does not depend on anyone's perceiving anything. We argue that the thesis survives the test when it is stated as a thesis about the founding of evidence. All thought is founded on perception, in the sense that its evidence depends, genetically and epistemically, on perceptual acts; formal reasoning is no exception, because its evidence rests ultimately on the perceptual acts of reading and writing in which derivations are produced and checked. The thesis does not make mathematical truth perceptual: what is founded is our access to the truths, not the truths themselves, which may be truths about abstract objects. We show that the thesis applies to itself without self-refutation, since its own evidence is founded on the reading of the arguments that support it, and we answer objections from anti-psychologism, mental calculation, the myth of the given and unsurveyed formal proofs.
\end{abstract}

\section{The problem}

In the lecture \enquote{The primacy of perception} and in the \emph{Phenomenology of Perception}, Merleau-Ponty argued that the perceived world is the always presupposed foundation of all rationality, all value and all existence \citep{MerleauPonty1945,MerleauPonty1964}. Perception is not one faculty among others, delivering data that intellect then interprets; it is the opening onto a world within which all intellectual operations take place, and the evidence of those operations is borrowed from it. The thesis continues Husserl's doctrine of founding. In the sixth \emph{Logical Investigation}, Husserl argued that categorial intuition, the intuition of states of affairs, collections and relations, is founded on sensuous intuition: I can see that the paper is white only on the basis of seeing the white paper \citep[VI, §§40--48]{Husserl1901}. In \emph{Experience and Judgment}, the founding is extended genetically: predicative judgment arises from prepredicative experience, and the forms of logic have their origin in the structures of perceptual explication \citep{Husserl1939}.

The hardest case for the thesis is formal reasoning. A logician deriving a theorem in a formal system does not perceive the objects the theorem is about, if it is about any objects; the validity of the derivation does not depend on anyone's perceiving it; and the truths of mathematics, on the view defended in this collection, are discovered truths about abstract objects (MTH-01, MTH-03). If formal reasoning is founded on perception, the founding must be of a kind that leaves these features intact. Merleau-Ponty himself said little about the case. His remarks on geometry, in the chapter on the cogito, argue that the evidence of a geometrical proof depends on the perceived figure and the motor activity of drawing it \citep[III.1]{MerleauPonty1945}; but the argument concerns Euclidean geometry, where the figure plays an obvious role, and leaves formal logic and algebra aside.

We argue that the thesis is true of formal reasoning as well, and that the way it is true is instructive. The evidence of a derivation rests on the perceptual acts of reading and writing in which the derivation is produced, inspected and checked. The founding is a founding of evidence, not of truth, and it is therefore compatible with any view of what mathematical truths are about. We then apply the thesis to itself.

\section{Framework}

\begin{definition}\label{d:founding}
An act $A$ is \emph{founded} on acts of kind $K$ if the evidence of $A$, that is, the fulfilment by which $A$ is justified for its subject, is possible only on the basis of acts of kind $K$. Founding is \emph{genetic} if acts of kind $K$ are required for the subject to have acquired the capacity for $A$, and \emph{epistemic} if they are required for each exercise of $A$ to have its evidence.
\end{definition}

Definition~\ref{d:founding} is deliberately restricted to evidence. That $A$ is founded on $K$-acts does not mean that the object of $A$ is a $K$-object, nor that what makes $A$ correct is anything about $K$-acts. Husserl's founding of categorial on sensuous intuition already has this form: the state of affairs that the paper is white is not itself sensuously perceived, though its intuition is founded on a sensuous perception of the paper \citep[VI, §48]{Husserl1901}.

\begin{definition}\label{d:perception}
A \emph{perceptual act} is an act that presents its object as bodily present, in any sensory modality: seeing, hearing, touching, and the kinaesthetic awareness of one's own movements. Reading and writing are perceptual acts in this sense: reading is the visual or tactile perception of inscriptions as inscriptions, and writing is the production of inscriptions under kinaesthetic and perceptual control.
\end{definition}

\begin{principle}\label{p:categorial}
Every categorial act is founded, epistemically and genetically, on sensuous or imaginative acts that present the material it articulates.
\end{principle}

\begin{principle}\label{p:imagination}
Imaginative acts are founded genetically on perceptual acts: the capacity to imagine an object of a given kind is acquired through the perception of objects of that kind.
\end{principle}

Principle~\ref{p:categorial} is Husserl's thesis in the sixth \emph{Investigation}. Principle~\ref{p:imagination} is his thesis that phantasy is a modification of perception, which presents its object \enquote{as if} it were present and draws its sensuous material from past perception \citep[§§39--40]{Husserl1939}.

\begin{assumption}\label{a:formal}
Formal reasoning consists of categorial acts whose material is presented by inscriptions: written formulae, diagrams and the configurations of signs in which a derivation is laid out, whether perceived or imagined.
\end{assumption}

Assumption~\ref{a:formal} is the claim that formal reasoning is reasoning in a notation. It does not say that formulae are what formal reasoning is about: a formula about numbers is about numbers. It says that the reasoner gains evidence that one formula follows from others by inspecting the formulae, that is, by recognizing their form, which is a categorial articulation of the perceived or imagined inscription. Historical work on the development of deduction supports the assumption: \citet{Netz1999} showed how Greek deductive practice was shaped by the lettered diagram and the formulaic language of the written proof, and the evidence of an algebraic computation is, as anyone who has made an error of transcription knows, inseparable from the correct reading of the lines on which it is performed.

\section{The argument}

\begin{lemma}\label{l:formal}
The evidence of formal reasoning is founded, epistemically, on perceptual or imaginative acts presenting inscriptions.
\end{lemma}

\begin{proof}
By Assumption~\ref{a:formal}, formal reasoning consists of categorial acts whose material is presented by inscriptions. By Principle~\ref{p:categorial}, each such act has its evidence only on the basis of the acts that present that material, which are perceptual acts of reading or imaginative acts presenting inscriptions.
\end{proof}

\begin{theorem}\label{t:primacy}
Formal reasoning is founded on perception: its evidence rests ultimately on perceptual acts of reading and writing.
\end{theorem}

\begin{proof}
By Lemma~\ref{l:formal}, the evidence of formal reasoning is founded on perceptual or imaginative presentations of inscriptions. Where the presentation is perceptual, it is reading, of an inscription that has been written. Where it is imaginative, Principle~\ref{p:imagination} founds it genetically on prior perception of inscriptions, which was again reading and writing. In either case the founding terminates in perceptual acts of reading and writing.
\end{proof}

\begin{corollary}\label{c:all}
All thought is founded on perception.
\end{corollary}

\begin{proof}
Formal reasoning is the case least obviously founded on perception, since its objects are not perceptual. Other thought is either perceptual, or categorially articulates perceived or imagined material and so is founded on perception by Principles~\ref{p:categorial} and~\ref{p:imagination}, or is reasoning, which proceeds in language, spoken or written, and so in perceived or imagined signs. Theorem~\ref{t:primacy} extends the founding to formal reasoning.
\end{proof}

\begin{proposition}\label{p:truth}
The founding established in Theorem~\ref{t:primacy} is a founding of evidence and does not make the truth of mathematical propositions depend on perception.
\end{proposition}

\begin{proof}
By Definition~\ref{d:founding}, the founding relation holds between acts and concerns their evidence. What makes a mathematical proposition true is a matter of what it is about, and the founding of the act in which it is known says nothing about that. A proposition about numbers is founded, as known by us, on the perception of numerals; it remains a proposition about numbers.
\end{proof}

Proposition~\ref{p:truth} shows how the primacy of perception coheres with the philosophy of mathematics defended in this collection. On MTH-01, numbers are abstract objects; on MTH-03, mathematical truths are discovered, while notations and axiomatizations are invented; on MTH-14, numbers are positions in structures. None of these theses is threatened by a thesis about the founding of evidence. MTH-02 holds that abstract objects are known through reliable belief-forming methods, and Theorem~\ref{t:primacy} says what those methods are like from the inside: they are practices of reading and writing in which categorial evidence is gained by inspecting perceived inscriptions. The two theses describe the same knowledge from different sides, the one reliabilist and the other phenomenological. And because the founding concerns evidence, not subject matter, Theorem~\ref{t:primacy} is not a version of formalism, which MTH-15 rejects. The formalist takes the inscriptions to be what mathematics is about; we take them to be what mathematical evidence is gained through.

The thesis also explains a feature of mathematical practice that the abstract-object view alone leaves unexplained: the importance of notation. Good notation makes proofs evident that bad notation hides, though the truths proved are the same. On the present account, this is because notation determines the perceptual material on which categorial evidence is founded, and some material lends itself to articulation better than other material. \citet{Giaquinto2007} documents the role of visual thinking in mathematics, and \citet{Parsons2008} argues that intuition of strings of symbols, as quasi-concrete objects, plays an ineliminable role in our knowledge of arithmetic. Both are consistent with, and partly explained by, Theorem~\ref{t:primacy}. Hilbert's finitism, which took the evidence of elementary arithmetic to rest on the surveyability of concrete signs \citep{Hilbert1926}, is a precursor, though Hilbert drew from it conclusions about the content of mathematics that we reject.

\section{Application to itself}

The thesis is a philosophical thesis, defended by reasoning. If all thought is founded on perception, so is the thought that all thought is founded on perception. Does the thesis then undermine itself?

\begin{proposition}\label{p:self}
The evidence of Corollary~\ref{c:all} is founded on perception, and this does not undermine it.
\end{proposition}

\begin{proof}
The evidence of Corollary~\ref{c:all}, for any reader, is gained by reading arguments such as the present one, and by reflection on one's own acts, which proceeds in language and articulates the remembered or presently lived acts reflected on. By Corollary~\ref{c:all} applied to these acts, the evidence is founded on perception. A thesis would be undermined by its self-application only if its application showed its evidence to be defective, and the thesis says nothing of the kind: it says that evidence is founded on perception, not that founded evidence is deficient. Founded evidence is evidence; the categorial evidence that the paper is white is not impaired by being founded on seeing the paper.
\end{proof}

The self-application is not idle. It shows that the reflective evidence on which phenomenology relies has the same structure as the evidence it describes, so that phenomenology does not claim for itself an access that it denies to other thought. This is the condition that, according to PHE-18, phenomenology must meet if it is to be a rigorous science in a sense it can define without circularity of the vicious kind. Merleau-Ponty saw the point when he described reflection as \enquote{radical} only if it takes account of its own dependence on an unreflective life which is its starting point \citep[Preface]{MerleauPonty1945}.

The self-application also has a historical dimension. Husserl argued, in the \emph{Origin of Geometry}, that ideal objects acquire their persisting objectivity through writing, which allows their sense to be transmitted and reactivated by later subjects \citep{Husserl1939b}. Writing is thus not only the medium in which formal evidence is produced but the condition under which it can be shared and checked by a community. That condition is perceptual through and through. PHE-17 examines the other side of the same fact: sedimented sense can be transmitted without being reactivated, and a science can come to operate with written formulae whose founding it has forgotten.

\section{Objections and replies}

\begin{objection}[Anti-psychologism]
Frege and the Husserl of the \emph{Prolegomena} argued that logic is independent of psychology: the laws of logic are not laws of thinking, and their validity does not depend on any mental act \citep{Frege1884}; \citep[Prolegomena]{Husserl1900}. To found logic on perception is a relapse into psychologism.
\end{objection}

\begin{reply}
Anti-psychologism concerns the validity of logical laws and the content of logical truths, and on both counts Proposition~\ref{p:truth} concedes everything. What is founded on perception is our evidence for logical truths, which is an act of a subject and has whatever structure acts of that kind have. Husserl did not see a conflict between the \emph{Prolegomena} and the sixth \emph{Investigation}, in which the founding of categorial intuition is established; the conflict would arise only if the founding of evidence were taken for the grounding of truth, which Definition~\ref{d:founding} excludes.
\end{reply}

\begin{objection}[Mental calculation]
A skilled calculator multiplies large numbers in her head without reading or writing anything. Her reasoning is formal, and its evidence does not rest on reading or writing.
\end{objection}

\begin{reply}
The calculator imagines inscriptions, or their auditory counterparts, and operates on them: experienced calculators report working with visualized digits or with recited number words. By Principle~\ref{p:imagination} these imaginative acts are founded genetically on the perception of written or spoken numerals, which is why the capacity is acquired by practice with written calculation. The founding in this case is genetic rather than epistemic, and Theorem~\ref{t:primacy} says only that the evidence rests \emph{ultimately} on reading and writing. The same holds of the logician who checks a short derivation without writing it down.
\end{reply}

\begin{objection}[The myth of the given]
\citet{Sellars1956} argued that no perceptual state has epistemic authority independently of conceptual capacities. If perception is itself conceptually informed, it cannot be the foundation of thought; the founding runs in both directions.
\end{objection}

\begin{reply}
Sellars's target is a foundationalism in which non-conceptual episodes justify beliefs by themselves. The founding of Definition~\ref{d:founding} is not of this kind. It does not say that perception justifies formal reasoning; it says that formal reasoning has its evidence only on the basis of perception. This is compatible with perception's being shaped by acquired concepts, as reading obviously is: one perceives a formula as a formula only after learning the notation. The primacy at issue is a primacy of perception within experience, not of a pre-conceptual given, and it is compatible with PHE-02's thesis that givenness is always givenness of something to someone as something.
\end{reply}

\begin{objection}[Unsurveyed proofs]
On MTH-09, a formally verified derivation is a proof whether or not anyone surveys it. Such a proof has evidence that does not rest on anyone's reading it.
\end{objection}

\begin{reply}
MTH-09 concerns what a proof is, and we accept it: proofhood does not depend on survey. Our thesis concerns the evidence that a subject has for a proposition proved in this way. The mathematician who accepts a machine-verified theorem has read its statement, has read the output of the checker and has, directly or through the community, read the code of the checker and the formalization of the definitions. Her evidence is founded on these acts of reading. It is a different evidence from that of a surveyed proof, mediated by trust in a mechanism rather than by inspection of each step, but it is founded on perception no less.
\end{reply}

\begin{objection}[The blind mathematician]
Many mathematicians have been blind, some from birth. If formal evidence rested on reading and writing, they would have none.
\end{objection}

\begin{reply}
Definition~\ref{d:perception} includes touch, hearing and kinaesthesis. Braille is read, dictated formulae are heard, and blind mathematicians typically report a rich imaginative manipulation of spatial and symbolic configurations founded on tactile and auditory experience. The thesis concerns perception, not vision, and the case confirms rather than refutes it: the blind mathematician's evidence rests on perceptual acts of a different modality.
\end{reply}

\section{What the result does not establish}

The result does not establish that mathematical objects are perceptual, constructed or mind-dependent; Proposition~\ref{p:truth} excludes this inference. It does not establish that formal reasoning is reducible to perception, or that its evidence is perceptual evidence: categorial evidence is founded on perceptual evidence, not identical with it. Finally, it does not settle how far the founding extends into concept acquisition in general. We have argued for the case of formal reasoning, which is the hardest, and for thought in general through Principles~\ref{p:categorial} and~\ref{p:imagination}; the reader who rejects the second principle will retain the epistemic founding of formal reasoning on present reading, and lose only its genetic founding in cases of purely imaginative reasoning. Recent phenomenological work on mathematical intuition \citep{Tieszen1989,Hopkins2011} develops related views in more detail.

\section{Conclusion}

All thought is founded on perception, including formal reasoning: the evidence of a derivation rests on the categorial articulation of perceived or imagined inscriptions, and imagined inscriptions are founded on the reading and writing of perceived ones. The founding is a founding of evidence, not of truth, and so leaves intact the abstract objects and discovered truths of mathematics. Applied to itself, the thesis shows that the evidence of phenomenology has the same structure as the evidence it describes, which is a condition of its coherence, not a threat to it.

\begin{thebibliography}{99}\small
\bibitem[Frege(1884)]{Frege1884} Frege, G. (1884). \emph{Die Grundlagen der Arithmetik}. Koebner.
\bibitem[Giaquinto(2007)]{Giaquinto2007} Giaquinto, M. (2007). \emph{Visual Thinking in Mathematics: An Epistemological Study}. Oxford University Press.
\bibitem[Hilbert(1926)]{Hilbert1926} Hilbert, D. (1926). Über das Unendliche. \emph{Mathematische Annalen}, 95, 161--190.
\bibitem[Hopkins(2011)]{Hopkins2011} Hopkins, B.~C. (2011). \emph{The Origin of the Logic of Symbolic Mathematics: Edmund Husserl and Jacob Klein}. Indiana University Press.
\bibitem[Husserl(1900)]{Husserl1900} Husserl, E. (1900). \emph{Logische Untersuchungen. Erster Teil: Prolegomena zur reinen Logik}. Niemeyer.
\bibitem[Husserl(1901)]{Husserl1901} Husserl, E. (1901). \emph{Logische Untersuchungen. Zweiter Teil: Untersuchungen zur Phänomenologie und Theorie der Erkenntnis}. Niemeyer.
\bibitem[Husserl(1939a)]{Husserl1939} Husserl, E. (1939a). \emph{Erfahrung und Urteil. Untersuchungen zur Genealogie der Logik}, ed. L.~Landgrebe. Academia.
\bibitem[Husserl(1939b)]{Husserl1939b} Husserl, E. (1939b). Die Frage nach dem Ursprung der Geometrie als intentional-historisches Problem, ed. E.~Fink. \emph{Revue internationale de philosophie}, 1(2), 203--225.
\bibitem[Merleau-Ponty(1945)]{MerleauPonty1945} Merleau-Ponty, M. (1945). \emph{Phénoménologie de la perception}. Gallimard.
\bibitem[Merleau-Ponty(1964)]{MerleauPonty1964} Merleau-Ponty, M. (1964). \emph{The Primacy of Perception}, ed. J.~M. Edie. Northwestern University Press.
\bibitem[Netz(1999)]{Netz1999} Netz, R. (1999). \emph{The Shaping of Deduction in Greek Mathematics: A Study in Cognitive History}. Cambridge University Press.
\bibitem[Parsons(2008)]{Parsons2008} Parsons, C. (2008). \emph{Mathematical Thought and Its Objects}. Cambridge University Press.
\bibitem[Sellars(1956)]{Sellars1956} Sellars, W. (1956). Empiricism and the philosophy of mind. In H.~Feigl and M.~Scriven (Eds.), \emph{Minnesota Studies in the Philosophy of Science}, Vol.~1, 253--329. University of Minnesota Press.
\bibitem[Tieszen(1989)]{Tieszen1989} Tieszen, R. (1989). \emph{Mathematical Intuition: Phenomenology and Mathematical Knowledge}. Kluwer.
\end{thebibliography}

\end{document}
